\section*{Chapitre \ref{ch:hf:lead-lag-and-cross-venue-trading} --- Avance-retard et négociation inter-plateformes}

\begin{solution}{exo:hf:lead-lag-and-cross-venue-trading:1}
Le marché~2 corrige (grand $\alpha_2$), donc le marché~1 mène~; $\gamma_1=0.18/(0.18+0.02)=0.9$.
\end{solution}

\begin{solution}{exo:hf:lead-lag-and-cross-venue-trading:2}
$67.5$ transactions par séance $\times0.363\times100\times\$0.01=\$24.50$.
\end{solution}

\begin{solution}{exo:hf:lead-lag-and-cross-venue-trading:3}
$97/7\approx14$.
\end{solution}

\begin{solution}{exo:hf:lead-lag-and-cross-venue-trading:4}
Lorsque l'avance est plus courte que le pas de la grille, les réactions des deux marchés à une même nouvelle tombent dans le même intervalle~: leurs innovations sont
corrélées et l'ordre de Cholesky détermine qui reçoit la part commune. Lorsque l'avance couvre plusieurs pas, l'innovation du meneur vient en premier et
celle du suiveur relève surtout de la correction d'erreur~; les innovations sont alors presque non corrélées et les bornes se rejoignent.
\end{solution}

\begin{solution}{exo:hf:lead-lag-and-cross-venue-trading:5}
Une fois que le suiveur a bougé, la transaction se fait au nouveau prix, sans écart à combler~: le mouvement attendu est nul et le preneur paie tout de même la
moitié de la fourchette et les frais, soit $-0.6$ tick, proche des $-0.54$ mesurés à une seconde.
\end{solution}

\begin{solution}{exo:hf:lead-lag-and-cross-venue-trading:6}
Entre 300 millisecondes ($+0.24$) et 500 ($-0.22$).
\end{solution}

\begin{solution}{exo:hf:lead-lag-and-cross-venue-trading:7}
0.444 tick par action à une milliseconde et 0.142 à 100~: avec une avance de 0.2 seconde, 100 millisecondes absorbent déjà les deux tiers de l'edge.
\end{solution}

\begin{solution}{exo:hf:lead-lag-and-cross-venue-trading:8}
La part dépend du pas d'échantillonnage~: sur des données à une seconde, une avance de quelques millisecondes est invisible et les parts reflètent surtout lequel des marchés a le bruit le plus faible
Le fonds peut compter aux fréquences où le desk négocie~; il faut mesurer sur ses propres ticks et à son propre horizon.
\end{solution}

\begin{solution}{pb:hf:lead-lag-and-cross-venue-trading:1}
\begin{enumerate}
\item L'essentiel de la formation des prix a lieu sur l'E-mini pour le S\&P 500 et le Nasdaq-100~; elle se partage entre le future standard et le fonds pour le S\&P 400~; le fonds S\&P 500
mène bien plus les fonds sectoriels que l'inverse.
\item Les contrats à terme ont porté l'essentiel de la formation des prix pour les obligations à cinq et dix ans à partir de 1999.
\item Corrélation médiane de 0.10 à 10 ms et de 0.008 à 1 ms en 2011.
\item Voir la \cref{def:hf:lead-lag-and-cross-venue-trading:lead}.
\item $\Delta p_t=\alpha(p^1_{t-1}-p^2_{t-1})+\sum\Gamma_i\Delta p_{t-i}+e_t$~; $\gamma_1=\alpha_2/(\alpha_2-\alpha_1)$~; la part d'information découle du
facteur de Cholesky, dans les deux ordres.
\item Parts d'information de 0.52--0.73, 0.81--0.90, 0.97--0.98 et 0.998 pour des avances de 0.05, 0.2, 0.5 et 1 seconde~; part de composante de 0.57, 0.73, 0.87,
0.98~; les bornes se resserrent à mesure que la grille résout l'avance.
\item L'échantillonnage sur une grille ajoute le biais d'Epps et brouille les avances plus courtes que le pas.
\item 0.05 et 0.25~; 0.225 et 0.25~; 0.475 et 0.575~; 1.075 et 1.1 secondes.
\item Les carnets d'ordres par défaut suivaient leur prix efficient sur quelques secondes, de sorte que le pic de corrélation croisée était plat~; les estimations étaient de $-1.65$ et $+1.0$
secondes.
\item Croiser la fourchette du suiveur pour un lot lorsque le meneur a bougé d'au moins un tick de plus sur la dernière seconde~; valorisation au mid du suiveur cinq
secondes plus tard, après la moitié de la fourchette et des frais de 0.1 tick.
\item 0.485, 0.43, 0.24, $-0.22$ et $-0.54$ tick par action.
\item L'avance injectée est d'une demi-seconde~; une avance réelle entre contrat à terme et fonds était de quelques millisecondes en 2011, de sorte que la courbe serait comprimée dans ces valeurs.
\item Il retire ou décale son côté périmé lorsque le meneur bouge (chapitre~7).
\item Injecté 0.5 s, estimé 0.475 et 0.575 s~; part d'information du meneur de 0.97 à 0.98~; l'edge disparaît entre 300 et 500 ms.
\item La liquidité, les horaires, les frais, les pannes et la concurrence font varier la formation des prix~; négocier l'avance elle-même la raccourcit.
\item Plateforme crypto menante vers la plateforme en retard~: Makarov et Schoar ont trouvé des écarts bien plus grands d'un pays à l'autre, là où les contrôles des capitaux ralentissent l'arbitrage.
\item Estimations glissantes de l'avance et des parts d'information par tranche horaire, avec des alertes lorsqu'elles sortent de leur plage habituelle.
\item Ce que le preneur gagne est le mouvement qu'il anticipe~; une fois plus rapide que l'avance, être encore plus rapide apporte peu sur ce marché (0.485 à une
milliseconde contre 0.463 à cinquante).
\item Rien n'est gagné avec la vitesse inutilisée tant que son coût fixe est payé (effet de levier opérationnel du chapitre~1).
\item La vitesse relative aux autres intervenants qui négocient la même avance.
\end{enumerate}
\end{solution}

\begin{solution}{iq:hf:lead-lag-and-cross-venue-trading:1}
Le mouvement est peut-être déjà intégré sur les autres plateformes du fonds~; les teneurs de marché du fonds ont pu se retirer et la fourchette s'élargir~; le mouvement du future
peut être une seule exécution qui revient sur ses pas~; l'écart peut ne pas couvrir la fourchette et les frais au moment où l'ordre arrive.
\iqlookfor{coûts, données périmées et risque d'arriver en retard.}
\end{solution}

\begin{solution}{iq:hf:lead-lag-and-cross-venue-trading:2}
L'échantillonnage sur grille apparie des accroissements sur des intervalles où une série n'a peut-être pas encore bougé~; une part du mouvement commun est donc manquée (l'effet Epps).
Hayashi--Yoshida somme les produits de chaque paire d'accroissements qui se chevauchent, sur les ticks propres à chaque série~; aucun mouvement n'est donc perdu.
\iqlookfor{la source du biais et l'idée de chevauchement.}
\end{solution}

\begin{solution}{iq:hf:lead-lag-and-cross-venue-trading:3}
Les innovations des deux marchés sont corrélées, et la décomposition de Cholesky attribue la part corrélée au marché placé en premier~;
les deux ordres encadrent la part. Elles sont proches lorsque les innovations sont presque non corrélées~: une grille fine, ou une avance plus longue que le pas.
\iqlookfor{l'argument de l'ordre.}
\end{solution}

\begin{solution}{iq:hf:lead-lag-and-cross-venue-trading:4}
Horodater au niveau de la carte réseau, avec des horodatages matériels disciplinés par une source de temps commune (precision time protocol vers une référence satellitaire), enregistrer
aussi les horodatages propres aux bourses et mesurer le délai de trajet entre les sites (livre~14, chapitre~4).
\iqlookfor{horodatages matériels et horloge partagée.}
\end{solution}

\begin{solution}{iq:hf:lead-lag-and-cross-venue-trading:5}
Quelqu'un de plus rapide négocie les mêmes écarts~: vos exécutions arrivent plus tard dans chaque opportunité. Vérifiez votre rang de latence (délais d'exécution comparés au mouvement du meneur), la
vitesse de retrait des cotations du suiveur et la durée propre de l'avance.
\iqlookfor{concurrence et façon de la mesurer.}
\end{solution}

\begin{solution}{iq:hf:lead-lag-and-cross-venue-trading:6}
$\alpha_\perp=(\alpha_2,0)$, donc $\gamma=(1,0)$. Lorsque le marché~1 est placé en premier, sa part d'information vaut un~; placé en second, elle vaut $1-\rho^2$, avec $\rho$
la corrélation des deux innovations~: les bornes sont $[1-\rho^2,1]$.
\iqlookfor{l'algèbre et le rôle de la corrélation des innovations.}
\end{solution}
